\section{Conclusion}
\label{sec:conclusion}

Phase~1 built a working DMM from textbook ingredients: an
Avellaneda--Stoikov core widened by a Guéant multi-level ladder, an
OFI skew, session-adaptive $k$, and a depth-gated hedge router on
B and C. With defaults calibrated by hand it already produced a
coherent, mean-reverting inventory trajectory and a positive
inception-spread P\&L on synthetic EUR/USD paths. Phase~2 replaced
those defaults with MLE estimates drawn from Phase~1's own fill
history --- arrival $k$, Pareto tail, GARCH volatility, and a
grid-optimised $\gamma$ --- and showed that the calibrated quoter
kept the same P\&L shape while halving quoted-spread volatility and
tightening inventory control. Phase~3 stressed the same calibrated
configuration against a latency-advantaged HFT on a scripted 15-day
calendar: fills per day fell 24\% and headline MtM compressed, but
the inception-spread leg ($+3{,}510$~USD) survived and the DMM
absorbed every one-sided window without firing a single hedge,
which is the behaviour the contract is actually paying for.

The approach has real limits. Phase~1's quoting choices are
heuristic: inventory penalty, OFI weight and session multipliers
are defensible but not optimal, and Phase~2 only re-estimates a
subset of them (client-flow $A$ and $k$ stay user-set by design).
The HFT in Phase~3 is a single reactive agent on a hand-written
schedule rather than a learned adversary, and hedging is confined
to EUR/USD on B and C with no cross-pair routing. All results are
on simulated paths: the stock, the client flow and the HFT
behaviour are all our own models, so the numbers should be read as
internally consistent rather than externally validated.

Natural extensions follow directly. Replacing the HFT schedule
with an adversary that learns to withdraw optimally against our
quoting policy would turn Phase~3 into a proper game and expose
whether the fallback premium is robust to strategic play.
Extending calibration to live EUR/USD tick data would close the
loop between the MLE machinery and the real arrival process.
Cross-pair hedging --- routing EUR risk through EUR/GBP or
EUR/JPY when B and C depth is thin --- would relax the
single-venue-pair assumption that currently bounds the hedge
router. Finally, the contract-design question raised at the end
of Phase~3 (availability-based DMM compensation rather than
volume-based) is the natural bridge from this simulator to the
market-microstructure policy literature, and the clearest direction
in which this project would want to be taken next.
